\section{Artifact System}

Artifacts provide a mechanism for objects to reference additional data within a MIFX package.

An artifact describes the relationship between an object and an associated resource such as a toolpath file, geometry file, or auxiliary JSON payload.

Artifacts allow objects to remain lightweight while linking to external data stored elsewhere in the package.

\subsection{Artifact Structure}

Artifacts are defined as objects within an \texttt{artifacts} array.

Example:

\begin{verbatim}
{
  "role": "toolpath",
  "kind": "apt",
  "path": "cldata/op-2.apt",
  "present": true
}
\end{verbatim}

Each artifact contains the following fields:

\begin{itemize}
\item \textbf{role} — describes the semantic purpose of the artifact
\item \textbf{kind} — describes the format or data type of the artifact
\item \textbf{path} — relative file path to the referenced resource
\item \textbf{present} — indicates whether the referenced resource is included in the package
\end{itemize}

\subsection{Artifact Roles}

The \texttt{role} field defines how the referenced resource should be interpreted.

Examples of artifact roles include:

\begin{itemize}
\item \texttt{toolpath}
\item \texttt{tool}
\item \texttt{operation\_csys}
\item \texttt{setup\_geometry}
\item \texttt{tool\_profile}
\item \texttt{in\_process\_mesh}
\end{itemize}

The specification does not restrict artifact roles to a fixed list.

Software implementations should ignore roles they do not recognize.

\subsection{Artifact Kinds}

The \texttt{kind} field describes the data format of the referenced resource.

Examples include:

\begin{itemize}
\item \texttt{json}
\item \texttt{apt}
\item \texttt{stl}
\item \texttt{3mf}
\item \texttt{step}
\end{itemize}

Implementations may support additional formats depending on their capabilities.

\subsection{Relative Paths}

Artifact paths are specified using relative paths within the package.

This ensures that the entire MIFX archive remains self-contained and portable.

Example:

\begin{verbatim}
"path": "../../tools/main/2.json"
\end{verbatim}

\subsection{Optional Artifacts}

The \texttt{present} field indicates whether the referenced resource exists in the package.

This allows a MIFX file to declare resources that may be generated or supplied later by other systems.

Software implementations should handle missing artifacts gracefully.